\subsection{应用}\label{sec:20e}
我们沿用高维情形中的策略。为二重分歧覆盖 $\pi:X=X_m^1\to\CC$ 选取消失路径基，考察 $\F(\pi)$ 中相应的全 $\Ai$ 子范畴，将其记为 $\B^\to$。显式地，它的对象为 $(Z_0,\ldots,Z_m)$，态射空间为
\[
\hom_{\B^\to}(Z_i,Z_j)=\begin{cases}
\K\cdot c_{ji}\oplus\K\cdot d_{ji},&i<j,\\
\K\cdot e_{ii},&i=j,\\
0,&i>j,
\end{cases}
\]
其中所有次数均为零。$e_{ii}$ 是严格单位元，除此之外，有
\[
\begin{aligned}
\mu_{\B^\to}^2(c_{kj},c_{ji})&=c_{ki},&
\mu_{\B^\to}^2(d_{kj},d_{ji})&=d_{ki},\\
\mu_{\B^\to}^2(d_{kj},c_{ji})&=0,&
\mu_{\B^\to}^2(c_{kj},d_{ji})&=0.
\end{aligned}
\]
\SourcePage{319}
这与高维情形中的对应表达式\eref{20.5}存在差异，归根结底是因为 $S^0$ 的上同调环与高维球面的上同调环有不同的性质。取扭曲复形 $(C_0,\ldots,C_m)$，定义为
\begin{equation}\label{eq:20.17}
C_0=Z_0,\qquad C_j=\Cone(c_{j,j-1}\pm d_{j,j-1}),\quad j>0,
\end{equation}
其中对每个 $j$ 都可自由选择符号；这一歧义将在下面的\stref{笔记}{20.5}中讨论。令 $\overline{\A}$ 为 $\Tw\B^\to$ 中由这些对象构成的全 $\Ai$ 子范畴。容易算得 $H(\overline{\A})\cong\overline A$。
\StatementNumber{3}\begin{lemma}\label{st:20.3}
$\overline{\A}$ 不是形式的。
\end{lemma}
\begin{proof}
显然，只须证明 $m=2$ 的情形，因为其余所有 $\overline{\A}$ 都包含它作为全 $\Ai$ 子范畴。由于集合\eref{20.17}生成 $\Tw(\B^\to)$，$\overline{\A}$ 上与 $\B^\to$ 上的扭曲复形 $\Ai$ 范畴拟等价，进而它们的 Hochschild 上同调群同构。对于代数，这类结果可追溯到\bcite{71}；证明稍作调整便适用于 $\Ai$ 情形。直接计算得到
\[
\HH(\B^\to,\B^\to)=
\underbrace{\K}_{\text{次数为零}}\oplus
\underbrace{\K^2}_{\text{次数为一}}.
\]
这与\stref{引理}{20.2}中部分算出的 $\overline A$ 的 Hochschild 上同调不同，由此得到所需的非形式性结论。
\end{proof}

令 $\A$ 为 $\overline{\A}$ 中以 $(C_1,\ldots,C_m)$ 为对象的全 $\Ai$ 子范畴。由前面的讨论，它也不是形式的，至少当 $m\ge5$ 时如此。另一方面，将\stref{命题}{18.21}适当推广，见\stref{笔记}{20.5}，可知 $\A$ 拟同构于 $\F$。因此：
\StatementNumber{4}\begin{corollary}\label{st:20.4}
对任意 $m\ge5$，$\F=\F_m^1$ 都不是形式的。
\end{corollary}
\StatementNumber{5}\begin{note}\label{st:20.5}
我们在二维情形中使用了\stref{命题}{18.21}，此时 Lefschetz 纤维化退化为二重分歧覆盖。虽然原先的讨论并未涵盖这种情形，但很容易调整其证明，也可改用更初等的论证。不过，写出精确陈述时须稍加谨慎；\stref{笔记}{19.10}已简要提到这一点。态射 $c_{j,j-1}$ 与 $d_{j,j-1}$ 对应于 Lefschetz 顶针之间的交点。相应的陈述是，匹配圈为 $c_{j,j-1}\pm d_{j,j-1}$ 的映射锥，而符号的选择决定匹配圈所带的 Pin 结构。在我们的情形中，匹配圈是在 $X$ 上构成 $(A_m)$ 构型的曲线；原先假设它们带有非平凡 Pin 结构，因为我们曾应用\stref{推论}{17.17}证明它们是分裂生成元。所幸，这些曲线恰好在
\SourcePage{320}
$H_1(X;\Z/2)$ 中线性无关，因而任意两组 Pin 结构的选择，都可通过用 $X$ 上的实线丛作扭曲而相互转换；按\secref{12i}的语言，这就是膜自同构。在代数上，这反映为 $\B^\to$ 具有改变 $c_{j,j-1}\pm d_{j,j-1}$ 中符号的自同构。因此，\eref{20.17}中符号的选择无关紧要。
\end{note}
